\section*{Question 2}

\partquestion{(a)}For LDA, we classify an instance $x$ into class $k$ with the largest
posterior probability. Using Bayes' theorem:
\[P(Y = k | X = x) = p_k(x) =  \frac{\pi_k f_k(x)}{\sum_{l = 1}^{K} \pi_l f_l(x)},\]
and substituting $f_k(x)$, we obtain the following equation:
\[
P(Y=k\mid X=x)
=
\frac{
\pi_k
\left[
\frac{1}{(2\pi)^{p/2}|\Sigma|^{1/2}}
\exp\left(
-\frac{1}{2}
(x-\mu_k)^T
\Sigma^{-1}
(x-\mu_k)
\right)
\right]
}
{
\sum_{l=1}^{K}
\pi_l
\left[
\frac{1}{(2\pi)^{p/2}|\Sigma|^{1/2}}
\exp\left(
-\frac{1}{2}
(x-\mu_l)^T
\Sigma^{-1}
(x-\mu_l)
\right)
\right]
}.
\]
By removing the common terms present in both the numerator and denominator, we get
\[
P(Y=k\mid X=x)
=
\frac{
\pi_k
\exp\left(
-\frac{1}{2}
(x-\mu_k)^T
\Sigma^{-1}
(x-\mu_k)
\right)
}
{
\sum_{l=1}^{K}
\pi_l
\exp\left(
-\frac{1}{2}
(x-\mu_l)^T
\Sigma^{-1}
(x-\mu_l)
\right)
}.
\]
To classify a value $X = x$, we need to see which of $p_k(x)$ is largest. Only the numerator is
needed because the denominator is the same for all the classes. By taking the logs,
and removing terms that do not depend on k, we obtain the following Linear Discriminant function
(Score Function) \parencite{gareth2023introduction}:
\[\delta_k(x) = x^T\hat{\Sigma}^{-1}\hat{\mu}_k - \frac{1}{2}\hat{\mu}_k^T\hat{\Sigma}^{-1}\hat{\mu}_k + \log(\hat{\pi}_k)\]

\partquestion{(b)}The training data was used to fit an LDA model. The \emph{LinearDiscriminantAnalysis}
class from \emph{scikit-learn} was used.

\begin{enumerate}[label=\roman*.]
    \item The fitted model estimates were found as follows:
    
          \textbf{Estimated prior probabilities:}

          \begin{table}[h]
            \centering
            \caption{Estimated Prior Probabilities}
            \begin{tabular}{cc}
            \hline
            Class $k$ & $\hat{\pi}_k$ \\
            \hline
            1 & 0.2466 \\
            2 & 0.3474 \\
            3 & 0.4060 \\
            \hline
            \end{tabular}
            \label{tab:prior_prob}
          \end{table}

          \newpage

          \textbf{Estimated class mean vectors (centroids):}

          \begin{table}[h]
            \centering
            \caption{Estimated Class Mean Vectors}
            \begin{tabular}{ccc}
            \hline
            Class $k$ & $\hat{\mu}_{k,1}$ & $\hat{\mu}_{k,2}$ \\
            \hline
            1 & 99.8582 & 100.0497 \\
            2 & 109.5319 & 104.6407 \\
            3 & 104.5715 & 119.3876 \\
            \hline
            \end{tabular}
            \label{tab:class_mean_vectors}
          \end{table}

          \textbf{Estimated shared covariance matrix:}

          \[
            \hat{\Sigma}=
            \begin{bmatrix}
            26.08973240 & 11.76180644 \\
            11.76180644 & 26.56801875
            \end{bmatrix},
          \]

          and 

          \[
            \hat{\Sigma}^{-1}=
            \begin{bmatrix}
            0.04788649 & -0.02119961 \\
            -0.02119961 & 0.04702442
            \end{bmatrix}
          \]

    \item The estimated discriminant functions of the fitted model in terms of $X_1$ and $X_2$ was obtained
          by substituting the fitted model estimate parameters found in (b) i. into the score function obtained
          in (a):

          \begin{align*}
            \delta_1(X)
            &=
            \begin{bmatrix}
            X_1 & X_2
            \end{bmatrix}
            \begin{bmatrix}
            0.04788649 & -0.02119961 \\
            -0.02119961 & 0.04702442
            \end{bmatrix}
            \begin{bmatrix}
            99.8582 \\
            100.0497
            \end{bmatrix}
            \\[1em]
            &\quad
            -\frac{1}{2}
            \begin{bmatrix}
            99.8582 & 100.0497
            \end{bmatrix}
            \begin{bmatrix}
            0.04788649 & -0.02119961 \\
            -0.02119961 & 0.04702442
            \end{bmatrix}
            \begin{bmatrix}
            99.8582 \\
            100.0497
            \end{bmatrix}
            +\log(0.2466)
            \\[1em]
            &=
            \begin{bmatrix}
            X_1 & X_2
            \end{bmatrix}
            \begin{bmatrix}
            2.6540 \\
            2.5811
            \end{bmatrix}
            - 527.4675 + \log(0.2466)
            \\[1em]
            &\underrightarrow{=-263.0309 + 2.6540X_1 + 2.5811X_2}.
          \end{align*}
          \begin{align*}
            \delta_2(X)
            &=
            \begin{bmatrix}
            X_1 & X_2
            \end{bmatrix}
            \begin{bmatrix}
            0.04788649 & -0.02119961 \\
            -0.02119961 & 0.04702442
            \end{bmatrix}
            \begin{bmatrix}
            109.5319 \\
            104.6407
            \end{bmatrix}
            \\[1em]
            &\quad
            -\frac{1}{2}
            \begin{bmatrix}
            109.5319 & 104.6407
            \end{bmatrix}
            \begin{bmatrix}
            0.04788649 & -0.02119961 \\
            -0.02119961 & 0.04702442
            \end{bmatrix}
            \begin{bmatrix}
            109.5319 \\
            104.6407
            \end{bmatrix}
            +\log(0.3474)
            \\[1em]
            &=
            \begin{bmatrix}
            X_1 & X_2
            \end{bmatrix}
            \begin{bmatrix}
            3.0189 \\
            2.5919
            \end{bmatrix}
            - 300.9440 + \log(0.3474)
            \\[1em]
            &\underrightarrow{=-302.0016 + 3.0189X_1 + 2.5919X_2}.
          \end{align*}
          \begin{align*}
            \delta_3(X)
            &=
            \begin{bmatrix}
            X_1 & X_2
            \end{bmatrix}
            \begin{bmatrix}
            0.04788649 & -0.02119961 \\
            -0.02119961 & 0.04702442
            \end{bmatrix}
            \begin{bmatrix}
            104.5715 \\
            119.3876
            \end{bmatrix}
            \\[1em]
            &\quad
            -\frac{1}{2}
            \begin{bmatrix}
            104.5715 & 119.3876
            \end{bmatrix}
            \begin{bmatrix}
            0.04788649 & -0.02119961 \\
            -0.02119961 & 0.04702442
            \end{bmatrix}
            \begin{bmatrix}
            104.5715 \\
            119.3876
            \end{bmatrix}
            +\log(0.4060)
            \\[1em]
            &=
            \begin{bmatrix}
            X_1 & X_2
            \end{bmatrix}
            \begin{bmatrix}
            2.4702 \\
            3.3885
            \end{bmatrix}
            - 331.4283 + \log(0.4060)
            \\[1em]
            & \underrightarrow{= -332.3276 + 2.4702X_1 + 3.3885X_2}.
          \end{align*}

    \item Thereafter, the fitted model was used to classify the test data. Figure \ref{fig:cm_lda}
          illustrates the confusion matrix obtained:

          \begin{figure}[h]
            \centering
            \includegraphics[width=0.55\textwidth]{images/confusion_matrix_lda.png}
            \caption{Confusion Matrix of the LDA model}
            \label{fig:cm_lda}
          \end{figure}

          Using Figure \ref{fig:cm_lda}, the overall accuracy of the model was obtained:

          \[Accuracy = \frac{42 + 86 + 128}{290} \times 100 = 88.28 \%. \]

          Since this is a multiclass classification problem, the sensitivity and specificity
          was computed for each class by obtaining the TP, FN, FP and TN for each class. Sensitivity
          is TP / (TP + FN) and specificity is TN / (TN + FP). Table \ref{tab:lda_class_metrics} illustrates the
          results: 

          \begin{table}[h]
            \centering
            \caption{Class-wise Sensitivity and Specificity}
            \begin{tabular}{ccccccc}
            \hline
            Class & TP & FN & FP & TN & Sensitivity & Specificity \\
            \hline
            1 & 42  & 22 & 12 & 214 & 0.6563 & 0.9469 \\
            2 & 86  & 11 & 21 & 172 & 0.8866 & 0.8912 \\
            3 & 128 & 1  & 1  & 160 & 0.9922 & 0.9938 \\
            \hline
            \end{tabular}
            \label{tab:lda_class_metrics}
          \end{table}

          Using the results obtained from Table \ref{tab:lda_class_metrics}, the following mean values
          were obtained:
          \[\text{Mean Sensitivity} = 0.8450\]
          \[\text{Mean Specificity} = 0.9440\]

\end{enumerate}



\partquestion{(c)-(e)iii.}

For the starting means, standard deviations and mixing proportions,
the following logic was used \parencite{blomer2013simple}.

\textbf{Mixing Proportions:} The starting mixing proportions for cluster $k$ is the fraction
of data points assigned to cluster $k$ by K-means:
\[
w_k = \frac{|C_k|}{|X|}
\]
\textbf{Means:} The starting means are simply the centroids for each cluster determined by K-means:
\[
\mu_k = \frac{1}{|C_k|}
\sum_{x \in C_k} x
\]
\textbf{Standard Deviations:} The starting standard deviation is the spatial spread of points assigned
to K-means cluster $k$:
\[
\Sigma_k =
\frac{1}{|C_k|}
\sum_{x \in C_k}
(x - \mu_k)(x - \mu_k)^T
\]
\newpage

\partquestion{(c)}\textbf{K = 2 components:}

\begin{enumerate}[label=\roman*.]
    \item K-means clustering was fit with 2 components to the combined data using the \emph{Kmeans} class
          from \emph{scikit-learn}. The \emph{random\_state} was set to 42 for consistency. Table \ref{tab:kmeans_centroids_2}
          illustrates the estimated centroids:

          \begin{table}[h]
            \centering
            \caption{Estimated Centroids for K-means with 2 components}
            \begin{tabular}{ccc}
            \hline
            Cluster & Centroid $X_1$ & Centroid $X_2$ \\
            \hline
            1 & 104.5338 & 101.6268 \\
            2 & 105.9013 & 118.9034 \\
            \hline
            \end{tabular}
            \label{tab:kmeans_centroids_2}
          \end{table}

    \item Using the fitted K-means model with 2 components, the following classified data
          can be seen in Figure \ref{fig:kmeans_scatter_2}:

          \begin{figure}[h]
            \centering
            \includegraphics[width=0.75\textwidth]{images/K2/Kmeans.png}
            \caption{Clustered Data using K-means with 2 components}
            \label{fig:kmeans_scatter_2}
          \end{figure}

    \item Thereafter, a GMM model with 2 components was fitted using the \emph{GaussianMixture}
          class from \emph{scikit-learn}. The \emph{random\_state} was set to 42 for consistency.
          The starting means, standard deviations and mixing proportions were obtained from the
          K-means model fitted in (i) and used during initialisation of the \emph{GaussianMixture} class.
          The starting values can be seen in Table \ref{tab:gmm_starting_values_2}:

          \begin{table}[h]
            \centering
            \caption{Starting Values for the GMM with 2 components}
            \begin{tabular}{cccccc}
            \hline
            Component & Mean $X_1$ & Mean $X_2$ & SD $X_1$ & SD $X_2$ & Mixing Proportion \\
            \hline
            1 & 104.5338 & 101.6268 & 6.8624 & 4.9744 & 0.5359 \\
            2 & 105.9013 & 118.9034 & 5.5837 & 5.0067 & 0.4641 \\
            \hline
            \end{tabular}
            \label{tab:gmm_starting_values_2}
          \end{table}

          \newpage

          Using the starting values seen in Table \ref{tab:gmm_starting_values_2}, a GMM model
          was fit and the following estimated values, seen in Table \ref{tab:gmm_estimated_values_2}, were obtained:

          \begin{table}[h]
            \centering
            \caption{Estimated Values for the GMM with 2 components}
            \begin{tabular}{cccccc}
            \hline
            Component & Mean $X_1$ & Mean $X_2$ & SD $X_1$ & SD $X_2$ & Mixing Proportion \\
            \hline
            1 & 105.6017 & 102.5052 & 7.0389 & 5.6883 & 0.5780 \\
            2 & 104.5752 & 119.4258 & 5.1567 & 5.0024 & 0.4220 \\
            \hline
            \end{tabular}
            \label{tab:gmm_estimated_values_2}
          \end{table}     
          
          Thereafter, the fitted GMM model's information criteria, $AIC$ and $BIC$, were obtained:
            \[
            \begin{aligned}
            \text{AIC} &= 19694.2276,\\
            \text{BIC} &= 19752.3001.
            \end{aligned}
            \]

    \item Using the fitted GMM model with 2 components, the following classified data
          can be seen in Figure \ref{fig:gmm_scatter_2}:

          \begin{figure}[h]
            \centering
            \includegraphics[width=0.75\textwidth]{images/K2/gmm.png}
            \caption{Clustered Data using GMM with 2 components}
            \label{fig:gmm_scatter_2}
          \end{figure}

\end{enumerate}

\newpage

\partquestion{(d)}\textbf{K = 3 components:}

\begin{enumerate}[label=\roman*.]
    \item K-means clustering was fit with 3 components to the combined data using the \emph{Kmeans} class
          from \emph{scikit-learn}. The \emph{random\_state} was set to 42 for consistency. Table \ref{tab:kmeans_centroids_3}
          illustrates the estimated centroids:

          \begin{table}[h]
            \centering
            \caption{Estimated Centroids for K-means with 3 components}
            \begin{tabular}{ccc}
            \hline
            Cluster & Centroid $X_1$ & Centroid $X_2$ \\
            \hline
            1 & 99.7470  & 99.3085 \\
            2 & 104.8763 & 120.0094 \\
            3 & 111.0364 & 106.0356 \\
            \hline
            \end{tabular}
            \label{tab:kmeans_centroids_3}
          \end{table}

    \item Using the fitted K-means model with 3 components, the following classified data
          can be seen in Figure \ref{fig:kmeans_scatter_3}:

          \begin{figure}[h]
            \centering
            \includegraphics[width=0.75\textwidth]{images/K3/Kmeans.png}
            \caption{Clustered Data using K-means with 3 components}
            \label{fig:kmeans_scatter_3}
          \end{figure}

    \item Thereafter, a GMM model with 3 components was fitted using the \emph{GaussianMixture}
          class from \emph{scikit-learn}. The \emph{random\_state} was set to 42 for consistency.
          The starting means, standard deviations and mixing proportions were obtained from the
          K-means model fitted in (i) and used during initialisation of the \emph{GaussianMixture} class.
          The starting values can be seen in Table \ref{tab:gmm_starting_values_3}:

          \begin{table}[h]
            \centering
            \caption{Starting Values for the GMM with 3 components}
            \begin{tabular}{cccccc}
            \hline
            Component & Mean $X_1$ & Mean $X_2$ & SD $X_1$ & SD $X_2$ & Mixing Proportion \\
            \hline
            1 & 99.7470  & 99.3085  & 4.5787 & 4.9459 & 0.3000 \\
            2 & 104.8763 & 120.0094 & 4.9881 & 4.3945 & 0.4028 \\
            3 & 111.0364 & 106.0356 & 4.0027 & 4.0317 & 0.2972 \\
            \hline
            \end{tabular}
            \label{tab:gmm_starting_values_3}
          \end{table}

          \newpage

          Using the starting values seen in Table \ref{tab:gmm_starting_values_3}, a GMM model
          was fit and the following estimated values, seen in Table \ref{tab:gmm_estimated_values_3}, were obtained:

          \begin{table}[h]
            \centering
            \caption{Estimated Values for the GMM with 3 components}
            \begin{tabular}{cccccc}
            \hline
            Component & Mean $X_1$ & Mean $X_2$ & SD $X_1$ & SD $X_2$ & Mixing Proportion \\
            \hline
            1 & 100.5453 & 99.6885  & 5.3160 & 5.2141 & 0.2975 \\
            2 & 104.6707 & 119.6564 & 5.0813 & 4.7909 & 0.4129 \\
            3 & 110.6269 & 105.6018 & 4.5360 & 4.5226 & 0.2896 \\
            \hline
            \end{tabular}
            \label{tab:gmm_estimated_values_3}
          \end{table}     
          
          Thereafter, the fitted GMM model's information criteria, $AIC$ and $BIC$, were obtained:
            \[
            \begin{aligned}
            \text{AIC} &= 19675.2533,\\
            \text{BIC} &= 19765.0017.
            \end{aligned}
            \]

    \item Using the fitted GMM model with 3 components, the following classified data
          can be seen in Figure \ref{fig:gmm_scatter_3}:

          \begin{figure}[h]
            \centering
            \includegraphics[width=0.75\textwidth]{images/K3/gmm.png}
            \caption{Clustered Data using GMM with 3 components}
            \label{fig:gmm_scatter_3}
          \end{figure}

\end{enumerate}

\newpage

\partquestion{(e)}\textbf{K = 4 components:}

\begin{enumerate}[label=\roman*.]
    \item K-means clustering was fit with 4 components to the combined data using the \emph{Kmeans} class
          from \emph{scikit-learn}. The \emph{random\_state} was set to 42 for consistency. Table \ref{tab:kmeans_centroids_4}
          illustrates the estimated centroids:

          \begin{table}[h]
            \centering
            \caption{Estimated Centroids for K-means with 4 components}
            \begin{tabular}{ccc}
            \hline
            Cluster & Centroid $X_1$ & Centroid $X_2$ \\
            \hline
            1 & 111.2671 & 105.8966 \\
            2 & 107.8834 & 122.4818 \\
            3 & 99.8886  & 98.5331 \\
            4 & 100.7348 & 115.4825 \\
            \hline
            \end{tabular}
            \label{tab:kmeans_centroids_4}
          \end{table}

    \item Using the fitted K-means model with 4 components, the following classified data
          can be seen in Figure \ref{fig:kmeans_scatter_4}:

          \begin{figure}[h]
            \centering
            \includegraphics[width=0.75\textwidth]{images/K4/Kmeans.png}
            \caption{Clustered Data using K-means with 4 components}
            \label{fig:kmeans_scatter_4}
          \end{figure}

    \item Thereafter, a GMM model with 4 components was fitted using the \emph{GaussianMixture}
          class from \emph{scikit-learn}. The \emph{random\_state} was set to 42 for consistency.
          The starting means, standard deviations and mixing proportions were obtained from the
          K-means model fitted in (i) and used during initialisation of the \emph{GaussianMixture} class.
          The starting values can be seen in Table \ref{tab:gmm_starting_values_4}:

          \begin{table}[h]
            \centering
            \caption{Starting Values for the GMM with 4 components}
            \begin{tabular}{cccccc}
            \hline
            Component & Mean $X_1$ & Mean $X_2$ & SD $X_1$ & SD $X_2$ & Mixing Proportion \\
            \hline
            1 & 111.2671 & 105.8966 & 3.8095 & 3.9774 & 0.2869 \\
            2 & 107.8834 & 122.4818 & 3.8208 & 3.6457 & 0.2303 \\
            3 & 99.8886  & 98.5331  & 4.6618 & 4.3288 & 0.2772 \\
            4 & 100.7348 & 115.4825 & 3.4678 & 3.8975 & 0.2055 \\
            \hline
            \end{tabular}
            \label{tab:gmm_starting_values_4}
          \end{table}

          \newpage

          Using the starting values seen in Table \ref{tab:gmm_starting_values_4}, a GMM model
          was fit and the following estimated values, seen in Table \ref{tab:gmm_estimated_values_4}, were obtained:

          \begin{table}[h]
            \centering
            \caption{Estimated Values for the GMM with 4 components}
            \begin{tabular}{cccccc}
            \hline
            Component & Mean $X_1$ & Mean $X_2$ & SD $X_1$ & SD $X_2$ & Mixing Proportion \\
            \hline
            1 & 110.6256 & 105.6406 & 4.5479 & 4.4650 & 0.2866 \\
            2 & 107.3547 & 121.6554 & 4.1547 & 4.1872 & 0.2327 \\
            3 & 100.6293 & 99.4669  & 5.4124 & 5.0472 & 0.2932 \\
            4 & 101.2129 & 116.7777 & 3.9741 & 4.4752 & 0.1875 \\
            \hline
            \end{tabular}
            \label{tab:gmm_estimated_values_4}
          \end{table}     
          
          Thereafter, the fitted GMM model's information criteria, $AIC$ and $BIC$, were obtained:
            \[
            \begin{aligned}
            \text{AIC} &= 19687.3218,\\
            \text{BIC} &= 19808.7462.
            \end{aligned}
            \]

    \item Using the fitted GMM model with 4 components, the following classified data
          can be seen in Figure \ref{fig:gmm_scatter_4}:

          \begin{figure}[h]
            \centering
            \includegraphics[width=0.75\textwidth]{images/K4/gmm.png}
            \caption{Clustered Data using GMM with 4 components}
            \label{fig:gmm_scatter_4}
          \end{figure}

\end{enumerate}

\partquestion{(f)}To determine the best GMM model, a AIC and BIC curve was plotted
against the number of components (clusters) to find a potential 'elbow' (Point at which
there is balance between model complexity and how well the model fits the data):

\begin{figure}[h]
    \centering
    \includegraphics[width=0.55\textwidth]{images/elbow.png}
    \caption{AIC and BIC values for different GMM models}
    \label{fig:elbow}
\end{figure}

According to Figure \ref{fig:elbow}, although the 2-component model has a slightly 
lower BIC, the 3-component model has the lowest AIC. Hence, the 3-component model
provides a more meaningful representation of the observed structure of the data. The 2-component
model does not fit the data well enough since it merges some groups together whereas the
4-component model is too complex, splitting groups unnecessarily \parencite{geron2019hands}. Thus, the 3-component
GMM model represents a model that provides the best balance of model complexity and model fit.

\partquestion{(g)}The GMM with 3 components was then used to classify the observations using
hard classification. The following confusion matrix was obtained:

\begin{figure}[h]
    \centering
    \includegraphics[width=0.55\textwidth]{images/confusion_matrix_gmm.png}
    \caption{Confusion Matrix of the GMM model with 3 components}
    \label{fig:cm_gmm}
\end{figure}

\newpage

From Figure \ref{fig:cm_gmm}, the following overall accuracy was obtained:

\[Accuracy = \frac{292 + 361 + 583}{1450} \times 100 = 85.24 \%. \]

The following class sensitivity and specificity was also obtained:

\begin{table}[h]
  \centering
  \caption{Class-wise Sensitivity and Specificity}
  \begin{tabular}{ccccccc}
  \hline
  Class & TP & FN & FP & TN & Sensitivity & Specificity \\
  \hline
  1 & 292 & 58  & 137 & 1013 & 0.8343 & 0.8809 \\
  2 & 361 & 139 & 56  & 944  & 0.7220 & 0.9440 \\
  3 & 583 & 17  & 21  & 879  & 0.9717 & 0.9767 \\
  \hline
  \end{tabular}
  \label{tab:gmm_class_metrics}
\end{table}

Using the results obtained from Table \ref{tab:gmm_class_metrics}, the following mean values
were obtained:
\[\text{Mean Sensitivity} = 0.8427\]
\[\text{Mean Specificity} = 0.9339\]

Comparing the results obtained to the ones obtained from the LDA model used in 2.(b), it is clear
that the LDA outperformed the GMM model. The sensitivity and specificity of the LDA model
seen in Table \ref{tab:lda_class_metrics} show an overall better model performance over the 
results seen in Table \ref{tab:gmm_class_metrics} for the GMM model with 3 compo The main reason has to do with the fact that the LDA
model was a supervised model which knew the group labels during training. The GMM model with 
3 components was able to identify group 3 very well but often confused true group 2 points as
group 1 instead. The model achieved an overall accuracy of 85.24\% which is not that far off
from the LDA's performance of 88.28\%. This shows that unsupervised learning can discover
meaningful insights from the data and pick up the latent structure oof the data even without
using true labels.

In conclusion, when class labels are available, a supervised approach is generally better
for accurate classification, but when class labels are not present, unsupervised approaches
such as Kmeans or GMMs, provide a useful way to identify natural clusters in the data. Hence,
unsupervised approaches can be very useful for tasks such as dimensionality reduction as well as
feature extraction.